

%%%%%%%%%%%%%%%%%%%%%%%%%%%%%%%%%%%%%%%%%%%%%%%%%%%%%%%%%%%%%
\subsubsection{具体分析}

问题一要求对零售数据集进行基本统计与多维度可视化分析，属于探索性数据分析（EDA）问题。本问采用描述性统计方法和多维分组聚合方法进行求解，为后续各问题的深入分析奠定数据认知基础。

首先对原始数据进行加载与质量检查，识别并处理缺失值和日期格式异常，接着利用pandas和numpy计算销售额的核心描述性统计量（均值、中位数、标准差、分位数、偏度、峰度等）以及总销售额、总订单数、有效客户数等核心业务指标，然后按日、月、年三级时间维度聚合销售额并绘制趋势折线图以直观展示整体走势与周期波动，最后按区域、产品大类、客户细分、发货模式等多维度分组统计销售额分布特征并绘制对比图表。具体分析流程如图\ref{fig:analysis1_flow}所示。

\begin{figure}[H]
  \begin{center}
    \begin{tikzpicture}[x=1cm, y=0.7cm]
      \node[box] (n11) at (0, 0) {数据加载解析};
      \node[box] (n12) at (5, 0) {缺失值处理};
      \node[box] (n13) at (10, 0) {描述性统计};
      \node[box] (n22) at (5, -3) {多维分组统计};
      \node[box] (n21) at (0, -3) {时间趋势分析};
      \node[box] (n23) at (10, -3) {业务指标计算};
      \node[box] (n31) at (0, -6) {趋势图绘制};
      \node[box] (n32) at (5, -6) {分布对比图};
      \node[box] (n33) at (10, -6) {结果总结};
      \draw[arrow] (n11) -- (n12);
      \draw[arrow] (n12) -- (n13);
      \draw[arrow] (n23) -- (n22);
      \draw[arrow] (n22) -- (n21);
      \draw[arrow] (n13) -- ++(0,-1.0) -| (n23);
      \draw[arrow] (n21) -- (n31);
      \draw[arrow] (n31) -- (n32);
      \draw[arrow] (n32) -- (n33);
    \end{tikzpicture}
    \caption{问题一分析流程图}
    \label{fig:analysis1_flow}
  \end{center}
\end{figure}

% ============================================================
\subsubsection{模型准备}

在进行统计分析和可视化之前，需要对原始数据进行必要的预处理。本节将从数据预处理和统计方法两方面进行准备。

原始数据来源于2022-2025年美国零售企业的全量订单数据集（data.csv），包含9800条订单明细记录和18个核心字段，覆盖订单基础信息、客户地理信息、产品分类信息和销售核心指标四大类信息。经检查发现以下问题：日期字段为DD/MM/YYYY格式需要专门解析；Postal Code字段存在11条缺失值（占比0.11\%）；数据编码为GBK格式需要转码处理；同一Order ID下可能包含多条Product ID的明细记录，订单维度分析需注意去重。

针对上述问题，按以下步骤进行预处理：第一步，使用GBK编码读取CSV文件，通过pandas的to\_datetime函数以dayfirst=True参数正确解析DD/MM/YYYY格式的日期字段。第二步，对Postal Code的11条缺失值，采用按所在州（State）分组取众数的方法进行填充，保持地理信息的一致性。第三步，从Order Date中提取年份、月份、季度、日期等时间衍生字段，扩展数据维度以支持多粒度的时间分析。第四步，针对描述性统计需要，对单笔销售额明细数据进行min/max/mean/median/std/skew/kurtosis等统计量计算；对订单维度去重统计时，使用nunique()方法对Order ID和Customer ID进行去重计数。

预处理后，数据规模保持不变（9800行×24列，新增6个时间衍生列），缺失值已全部处理完毕。销售额分布呈现高度右偏特征（偏度12.98，峰度304.45），绝大部分单笔交易集中在较低金额区间（中位数54.49美元），少数大额交易拉高了均值（230.77美元），这种长尾分布是零售行业的典型特征。

综上所述，数据预处理完成，为后续统计分析和可视化提供了清洁规整的输入数据。
